% Part: first-order-logic
% Chapter: lindstrom
% Section: introduction

\documentclass[../../../include/open-logic-section]{subfiles}

\begin{document}

\olfileid[tr]{mod}{lin}{int}
\section{Giriş}

Bu bölümde, birinci dereceden mantığın, bazı ek koşullar altında Tıkızlık
Teoremi ile Aşağı L\"owenheim--Skolem Teoremi'nin geçerli olduğu en geniş
mantık olarak Lindstr\"om tarafından verilen karakterizasyonunu ispatlamayı
amaçlıyoruz (\olref[fol][com][com]{thm:compactness} ve
\olref[fol][com][dls]{thm:downward-ls}). Önce teoremin uygulandığı genel
mantık sınıfını daha geniş bir çerçevede tanımlamamız gerekir. Kendimizi
\emph{bağıntısal} dillerle, yani yalnızca !!{predicate}s ve bireysel sabitler
içeren, fakat !!{function}s içermeyen dillerle sınırlayacağız.

\end{document}

